\section{Proofs}\label{app:proofs}
\paragraph{Invariance law.}
Under $B\mapsto BM$, $[A'_i,A'_j]=\sum_{k,l}M_{ki}M_{lj}[A_k,A_l]$, so the commutator matrix becomes
$C(M\otimes M)$, while $P_0(BM)=P_0(B)$. Hence $R\mapsto R(M\otimes M)$. The two-sided bound follows from
$\sigma_{\min/\max}(M\otimes M)=\sigma_{\min/\max}(M)^2$. For $\varepsilon>0$,
$P_\varepsilon(BM)=BM(M^\top GM+\varepsilon I)^{-1}M^\top B^\top=B(G+\varepsilon M^{-\top}M^{-1})^{-1}B^\top$.

\paragraph{Proposition~\ref{prop:integrate}.}
The commutator is bilinear and antisymmetric and satisfies the Jacobi identity in every associative algebra,
so closure makes $\mathfrak a$ a Lie subalgebra. The Lie subgroup theorem
\citep[Thm.~5.20]{hall2015lie} gives the unique connected Lie subgroup with Lie algebra $\mathfrak a$; it is
generated by $\exp(\mathfrak a)$. An irrational line in a torus shows that it need not be closed.
$\square$

\paragraph{Lemma~\ref{lem:delta}.}
Write $X=\sum_ix_iA_i$, $Y=\sum_jy_jA_j$ and $r_{ij}=N[A_i,A_j]$. By Cauchy--Schwarz over the $K^2$ index
pairs, $\|N[X,Y]\|=\|\sum x_iy_jr_{ij}\|\le\|x\|\|y\|(\sum\|r_{ij}\|^2)^{1/2}=\|x\|\|y\|K\sqrt{\Lc}$. Finally,
$\|X\|^2=x^\top Gx\ge\lambda_{\min}(G)\|x\|^2$, and likewise for $Y$. $\square$

\paragraph{Proposition~\ref{prop:bch}.}
Dynkin's form of the BCH series \citep{bonfiglioli2012topics} is
$\log(e^Xe^Y)=\sum_{m\ge1}\frac{(-1)^{m-1}}{m}\sum\frac{[w]}{|w|\prod_i r_i!s_i!}$. The inner sum runs over
words $w=X^{r_1}Y^{s_1}\cdots X^{r_m}Y^{s_m}$ with $r_i+s_i>0$; $[w]$ is the right-nested commutator and
$|w|=\sum_i(r_i+s_i)$.

Since the Frobenius norm is submultiplicative, $\|[w]\|\le2^{|w|-1}r^{|w|}$. Let $e_n$ bound $\|N[w]\|$ over
words of length $n$; then $e_1=0$ and $e_2\le\delta r^2$. Write $[w]=[W,Z]$ with $W\in\{X,Y\}$ and $Z$ a
nested commutator of length $n-1$. The term $[W,P_0Z]$ has both arguments in $\mathfrak a$, so its escape is
at most $\delta r\|Z\|\le\delta2^{n-2}r^n$, while $\|[W,NZ]\|\le2re_{n-1}$. By induction,
$e_n\le(n-1)2^{n-2}\delta r^n$.

Each word therefore contributes at most $\frac{\delta}{4}(2r)^{|w|}/\prod r_i!s_i!$. Summing over all tuples
gives $\sum_m\frac1m(e^{4r}-1)^m=-\ln(2-e^{4r})$ for $r<\ln2/4$. Subtracting the length-1 words, which lie in
$\mathfrak a$ and account for $4r$, gives the first bound.

The length-2 words sum exactly to $\frac12[X,Y]$, whose escape is at most $\frac12\delta r^2$; in the series
bound they account for $4\delta r^2$. Replacing that amount gives the second bound. Since $-\ln(2-e^u)-u$ has
non-negative Taylor coefficients, the ratio $(\varphi(r)-4r^2)/r^3$ is increasing, and $C(r_0)$ bounds it on
$[0,r_0]$.

Finally, the series converges absolutely on this ball, and $\|e^Xe^Y-I\|\le e^{2r}-1<1$ there. Both sides are
analytic on the connected ball and agree near $0$ \citep[Thm.~5.3]{hall2015lie}, so they agree on the whole
ball. $\square$

\paragraph{Proposition~\ref{prop:chain}.}
With $d_t=\hat z_t-z_t$ we have $d_{t+1}=R_td_t-e_t$ and $\|R_td_t\|=\|d_t\|$; iterate. $\square$

\paragraph{Proposition~\ref{prop:classify}.}
Subalgebras of $\so(d)$ are compact, hence reductive.

(a)~A 2-dimensional reductive Lie algebra is abelian. Since $\operatorname{rank}\so(5)=2$, it is a maximal
abelian subalgebra, i.e.\ a Cartan subalgebra, and all Cartan subalgebras of a compact Lie algebra are
conjugate \citep{adams1969lectures}.

(b)~A 3-dimensional subalgebra $\mathfrak h$ of $\so(4)$ is reductive and cannot be abelian, because the rank
is 2. Hence $\mathfrak h\cong\su(2)$. Its projections onto the two factors of
$\so(4)\cong\su(2)\oplus\su(2)$ are each either zero or an isomorphism. If one projection is zero,
$\mathfrak h$ is a factor. Otherwise $\mathfrak h$ is the graph of an automorphism, which is inner, so
$\mathfrak h$ is conjugate to the diagonal subalgebra. The diagonal subalgebra acts on
$\mathbb R^4\cong\mathbb C^2\otimes\mathbb C^2$ as $\mathbf 3\oplus\mathbf 1$, while a factor acts
isoclinically. The participation ratio is a conjugation invariant that separates the two classes ($2$ vs.\
$4$). A reflection in $O(4)$ exchanges the two factors.

(c)~A 4-dimensional subalgebra is reductive of rank at most 2, so it is $\su(2)\oplus\mathfrak u(1)$, with the
$\su(2)$ a factor (by the argument in (b)) and the $\mathfrak u(1)$ in the other factor. Its elements $(X,Y)$
have $Y$ confined to a line, whereas the diagonal $\so(3)$ consists of all $(X,X)$. Equivalently, the
centralizer of the diagonal $\so(3)$ in $\so(4)$ is trivial, so a fourth generator can never be added while
preserving closure. $\square$
